% ============================================================
% 1. INTRODUCTION
% ============================================================
\section{Introduction}
\label{sec:intro}

% --- Module 1: Why this problem matters (~45%) ---

Large language models (LLMs) have rapidly advanced across a wide
range of tasks, from question answering and code generation to web
navigation and embodied decision making~\cite{yao2023react,
shinn2023reflexion, zhou2024lats}. Yet, many tasks remain beyond the
reach of a single forward pass: complex reasoning, multi-step
planning, and open-ended exploration often require the model to
evaluate alternatives, verify intermediate results, or search over
candidate actions~\cite{yao2023tree, hao2023reasoning,
snell2024scaling}. This has motivated \emph{test-time compute}, where additional
computation at inference time improves decision quality. Across
multiple benchmarks, test-time optimization allows smaller models to
match or exceed the performance of larger ones at lower total
cost~\cite{snell2024scaling, wang2026flare}.

% While effective, test-time compute is expensive, incurring substantial per-step overhead (e.g., requiring multiple rollouts or candidate evaluations at each decision point)~\cite{yao2023tree, zhou2024lats}. As
% agents are deployed across diverse settings, applying this overhead
% uniformly becomes wasteful when only a fraction of steps benefit.
% This has motivated \emph{adaptive gating}: using uncertainty
% signals to decide \emph{when} extra compute is worthwhile.
% Prior work explores diverse mechanisms, including entropy thresholding~\cite{seag2025, cats2025, corefine2026}, calibrated confidence~\cite{s1_2025, han2025tokenbudget}, vote disagreement~\cite{catts2026}, and reinforcement learning~\cite{adaptthink2025, thinkless2025}. Across these mechanisms, adaptive gating consistently reduces cost while preserving or improving task performance within each method's evaluation setting.
While effective, test-time compute is expensive, incurring substantial per-step overhead (e.g., requiring multiple rollouts or candidate evaluations at each decision point)~\cite{yao2023tree, zhou2024lats}. As
agents are deployed across diverse settings, applying this overhead
uniformly becomes wasteful when only a fraction of steps benefit.
This has motivated \emph{adaptive gating}: using uncertainty
signals to decide \emph{when} extra compute is worthwhile.
Prior work explores diverse mechanisms, including entropy thresholding~\cite{seag2025, cats2025, corefine2026}, calibrated confidence~\cite{s1_2025, han2025tokenbudget}, vote disagreement~\cite{catts2026}, and learned difficulty estimators based on hidden states or reinforcement learning~\cite{diffadapt2025, zhu2025llmalreadyknows, adaptthink2025, thinkless2025}. Across these mechanisms, adaptive gating consistently reduces cost while preserving or improving task performance within each method's evaluation setting.

% Yet these gains often fail to transfer: methods calibrated for one task type frequently fail when applied to others~\cite{tao2025revisiting,
% heo2025llmuncertainty}. Vote-based
% gating designed for web agents degrades on fact-verification tasks;
% calibration-based methods that perform well on QA degrade in code
% generation. In several cases, adaptive gating \emph{actively
% harms} performance compared to never invoking the optimizer. This
% raises a central question: \emph{do uncertainty signals have a
% universal relationship with compute value, or does this relationship
% depend on the environment?}
However, this consensus rests on a narrow empirical base. The
fixed-direction assumption holds well in the settings where these
methods have been developed and evaluated, predominantly homogeneous
reasoning benchmarks where signal direction is consistent by
construction. Whether it transfers to heterogeneous agent
environments is an open question prior work has not directly
addressed, and we show that it does not: the relationship between
uncertainty and compute value can reverse direction across
environments, and adaptive gating can actively underperform never
invoking the optimizer. This raises a central question:
\emph{what determines the relationship between uncertainty signals
and the value of computation, and can this relationship be learned
in a way that generalizes across environments?}

% --- Module 2: Our finding + diagnosis (~30%) ---

% We conduct systematic measurement across eight agent environments spanning QA~\cite{yang2018hotpotqa}, code generation~\cite{hendrycks2021apps}, web navigation~\cite{yao2022webshop}, fact verification~\cite{thorne2018fever}, and text games~\cite{jansen2023twexpress}, using three model backbones. The signal–utility direction \emph{reverses} across environments: the same signal that predicts rollout benefit in code generation predicts rollout harm in fact verification, and can further reverse across model backbones within the same environment. Beyond direction, the most informative signal itself varies across environments, with structural features dominating in some while entropy carries no information.
We answer this through systematic measurement across six agent environments spanning QA~\cite{yang2018hotpotqa}, code
generation~\cite{hendrycks2021apps}, web navigation~\cite{yao2022webshop}, fact verification~\cite{thorne2018fever}, text games~\cite{jansen2023twexpress}, and crafting planning, using three model backbones. The signal--utility relationship reverses direction across environments:
the same signal that predicts rollout benefit in code generation predicts rollout harm in fact verification, and can further reverse
across model backbones within the same environment. Moreover, the most informative signal itself varies across environments, with
structural features dominating in some settings while entropy carries little or no information in others.
 
% We trace this reversal to a conflation of two uncertainty sources.
% \emph{Information poverty} arises when the agent lacks sufficient
% data; here rollouts from the same model cannot compensate, and high
% entropy predicts rollout harm. \emph{Decision difficulty} arises
% when the agent faces multiple viable paths; here rollouts explore
% alternatives productively, and high entropy predicts rollout benefit.
% Environments differ in the proportion of these two state types,
% causing the same signal to carry opposite meaning. This is an
% instance of Simpson's paradox~\cite{simpson1951interpretation,kdd2023simpson} in the signal--utility space. We prove
% that this has a sharp consequence: no fixed-direction gate can
% guarantee non-negative value of computation across environments with
% opposing directions (Proposition~\ref{prop:necessity}). Once the
% direction is wrong, more precise calibration makes performance
% \emph{worse}, because the gate more precisely targets harmful states.
We trace this reversal to a conflation of two distinct sources of uncertainty. \emph{Information poverty} (Type~I) arises when the
agent lacks sufficient information; rollouts from the same model cannot compensate, and high uncertainty predicts \emph{negative}
value of computation. In contrast, \emph{decision difficulty} (Type~D) arises when multiple viable options exist; additional
computation helps explore alternatives, and high uncertainty predicts \emph{positive} value. Code generation often involves
multiple valid implementations (Type~D), while fact verification may lack sufficient evidence altogether (Type~I). Because
environments differ in their mixture of these two state types, the same uncertainty signal carries opposite meanings depending on
which source dominates, a phenomenon closely related to Simpson's paradox~\cite{simpson1951interpretation,kdd2023simpson}. The
consequence is sharp: no fixed-direction gating strategy can avoid net-harmful gating across environments with opposing signal
directions (Proposition~\ref{prop:necessity}), and once the direction is misaligned, improved calibration can \emph{worsen} performance by more precisely selecting harmful states.

% --- Module 3: Method + contributions (~25%) ---

% Motivated by this analysis, we propose \textbf{EAAG}
% (\textbf{E}nvironment-\textbf{A}ware \textbf{A}daptive
% \textbf{G}ating), which discovers environment-specific gating
% patterns through randomized \emph{exploration}, LLM-based
% \emph{reasoning} about exploration data, and sparse
% \emph{learning} of signal direction and importance.
% Since direction is the primary determinant of gate quality, not
% model capacity, EAAG pairs a rich discovery process with a
% deliberately simple gate that adds zero per-step deployment cost.
% EAAG outperforms all fixed-direction baselines across eight
% environments and three LLM backbones.
Motivated by this analysis, we propose \textbf{EAAG} (\textbf{E}nvironment-\textbf{A}ware \textbf{A}daptive \textbf{G}ating), a deliberately minimal method that learns environment-specific signal–utility relationships from interaction data. Since the primary challenge is identifying signal direction rather than modeling capacity, EAAG's contribution is not a new gate architecture but a pipeline that reduces the adaptive-gating problem to a form where a standard sparse linear classifier suffices, adding zero per-step cost at deployment. Across six environments and three LLM backbones, EAAG outperforms all fixed-direction baselines.
%Importantly, since signal direction cannot be inferred from uncertainty alone and varies across environments, it must be discovered through interaction data, rather than assumed or calibrated from static signals.

Taken together, these results point to a more fundamental
observation: uncertainty signals like token entropy are not
sufficient statistics for the value of computation, because their
meaning depends on a latent state property that the signal alone
cannot identify. Direction discovery is the minimal mechanism that
closes this gap.

Our main contributions are as follows:
\begin{itemize}[leftmargin=*,nosep]
\item \textbf{A structural limitation of uncertainty-based gating.}
We show that uncertainty signals are not sufficient statistics for
the value of computation: their meaning depends on a latent state
type, and as a consequence no fixed-direction gate can be
universally correct.

\item \textbf{Empirical finding.} We demonstrate systematic reversal of signal--utility direction across environments and model backbones, and show that the most informative signals vary across settings.

\item \textbf{Adaptive gating method.} We propose EAAG, which learns signal direction per environment instead of assuming it, outperforming all fixed-direction baselines across six environments with zero deployment overhead.
\end{itemize}
